\section{الخاتمة}
\label{sec:conclusion}

يقدم هذا البحث إطارًا لإخفاء المعلومات الدلالي الواعي بالسياق في النقاشات المتسلسلة، يعيد تعريف مجال التضمين بوصفه \textit{قناة اختيار}. ومن خلال الانتقال من التلاعب على مستوى الرموز إلى ترميز بتات الحمولة في استهداف الرد واختيار الاستطرادة، تهدف الطريقة إلى الحفاظ على الطبيعية السياقية في النص المرئي مع الاحتفاظ بالتوافق مع واجهات نماذج اللغة الكبيرة وفق نموذج الصندوق الأسود.

يتضمن التنفيذ آلية تحقق مغلقة الحلقة تحاكي فك الترميز قبل النشر وتعيد التوليد تحت التوليد العشوائي. ويورد تقييمنا المبني على مخرجات التنفيذ مقاييس الطبيعية وتباعد أحاديات الكلمات على عينات ناجحة أكدها المستقبل، كما يصف أنماط الإخفاق المهيمنة عند استنفاد محاولات الإعادة. وسيركز العمل المستقبلي على تحسين معدلات النجاح من البداية إلى النهاية في سلاسل حقيقية جديدة، وتوسيع أحجام العينات النظيفة، وإجراء تقييمات مخصصة لكشف الإخفاء ضمن نماذج خصومية صريحة.
